\documentclass[12pt]{article}
\usepackage[utf8]{inputenc}
\usepackage[spanish]{babel}
\usepackage{amsmath}
\usepackage{amssymb}
\usepackage{booktabs}
\usepackage{geometry}
\geometry{margin=2.5cm}
\usepackage[most]{tcolorbox}
\newtcolorbox{intuicion}{
  colback=blue!5!white,
  colframe=blue!40!black,
  fonttitle=\bfseries,
  title=Intuici\'on,
  boxrule=0.6pt,
  arc=4pt,
  left=6pt, right=6pt, top=4pt, bottom=4pt
}
\newtcolorbox{intuicionmat}{
  colback=green!5!white,
  colframe=green!40!black,
  fonttitle=\bfseries,
  title=\textit{?`Qu\'e dice esta ecuaci\'on?},
  boxrule=0.6pt,
  arc=4pt,
  left=6pt, right=6pt, top=4pt, bottom=4pt
}
\newtcolorbox{explicacion}{
  colback=orange!5!white,
  colframe=orange!50!black,
  fonttitle=\bfseries,
  title=Explicaci\'on,
  boxrule=0.6pt,
  arc=4pt,
  left=6pt, right=6pt, top=4pt, bottom=4pt
}

\title{\textbf{Resumen: Cap\'itulo 3}\\
\large Modelos de Precios de T\'itulos en Mercados Financieros\\
\normalsize \textit{Mathematics for Finance: An Introduction to Financial Engineering}}
\author{}
\date{Junio 2026}

\begin{document}
\maketitle
\tableofcontents
\newpage

%---------------------------------------------------------------
\section{Modelos en Tiempo Continuo: Tasa Real de Inter\'es\\
\textit{(Continuous-Time Models: Real Interest Rate)}}
%---------------------------------------------------------------

\subsection{Din\'amica del Activo Libre de Riesgo \textit{(Risk-Free Asset Dynamics)}}

En tiempo continuo se asume que la \textbf{inversi\'on libre de riesgo (risk-free investment)} $B(t)$ sigue la din\'amica:

\begin{equation}
  dB(t) = B(t)\,r(t)\,dt
\end{equation}

\begin{intuicionmat}
El activo libre de riesgo crece a una tasa instant\'anea $r(t)$, que puede ser estoc\'astica. Al no tener componente de movimiento Browniano, $B(t)$ es \textbf{localmente predecible (locally predictable)}: conociendo $r(t)$, podemos aproximar el incremento $\Delta B(t)$ en un intervalo corto $\Delta t$ con certeza. La tasa $r$ es la \textbf{tasa nominal de inter\'es (nominal interest rate)}.
\end{intuicionmat}

\begin{intuicion}
Piensa en una cuenta bancaria que acumula intereses continuamente. Si la tasa es constante $r$, el saldo crece de forma exponencial y predecible: no hay sorpresas en el banco. Esto contrasta con las acciones, donde el precio puede saltar de forma inesperada.
\end{intuicion}

\subsection{Din\'amica del Nivel General de Precios \textit{(Price Level Dynamics)}}

Para introducir tasas reales, primero se modela el \textbf{nivel general de precios (general price level)} mediante el proceso $p(t)$:

\begin{equation}
  dp(t) = p(t)\bigl[\mu(t)\,dt + \sigma(t)\,dW(t)\bigr]
\end{equation}

\begin{intuicionmat}
El nivel de precios tiene dos componentes: una deriva determinista $\mu(t)\,dt$ (la inflaci\'on esperada) y un t\'ermino estoc\'astico $\sigma(t)\,dW(t)$ (choques inesperados de inflaci\'on). Cuando $\mu$ y $\sigma$ son constantes, este proceso es un \textbf{movimiento Browniano geom\'etrico (geometric Brownian motion)}.
\end{intuicionmat}

Bajo condiciones suaves sobre $\mu$ y $\sigma$, la soluci\'on expl\'icita del proceso --- verificable por la \textbf{regla de It\^o (It\^o's rule)} --- es:

\begin{equation}
  p(t) = p(0)\,\exp\!\left(\int_0^t\!\left[\mu(s) - \tfrac{1}{2}\sigma^2(s)\right]ds + \int_0^t\!\sigma(s)\,dW(s)\right)
\end{equation}

\begin{intuicionmat}
La f\'ormula expl\'icita de $p(t)$ dice que el nivel de precios es el valor inicial $p(0)$ multiplicado por un factor exponencial. La correcci\'on $-\frac{1}{2}\sigma^2$ en la integral determinista proviene de la regla de It\^o: es el precio de la no-linealidad al pasar de la din\'amica diferencial a la soluci\'on en niveles. A diferencia de $B(t)$, el proceso $p(t)$ \textbf{no} es localmente predecible por la presencia del movimiento Browniano.
\end{intuicionmat}

\begin{intuicion}
El nivel de precios se mueve como una acci\'on: tiene una tendencia promedio (la inflaci\'on esperada) pero tambi\'en recibe choques aleatorios que no podemos anticipar. Por eso la inflaci\'on es, en parte, impredecible.
\end{intuicion}

\subsection{Tasa Real de Inter\'es en Tiempo Continuo \textit{(Real Interest Rate in Continuous Time)}}

La \textbf{tasa real de inter\'es (real interest rate)} corresponde al retorno del activo libre de riesgo deflactado por el nivel de precios. El retorno relativo $d(B/p)/(B/p)$, calculado por la regla de It\^o, es:

\begin{equation}
  \frac{d\!\left(\dfrac{B(t)}{p(t)}\right)}{\dfrac{B(t)}{p(t)}} = \bigl[r(t) - \mu(t) + \sigma^2(t)\bigr]\,dt - \sigma(t)\,dW(t)
\end{equation}

\begin{intuicionmat}
Esta ecuaci\'on descompone el retorno real en dos partes: (1) un componente predecible $[r - \mu + \sigma^2]\,dt$, que es la tasa nominal menos la inflaci\'on esperada m\'as una correcci\'on cuadr\'atica por la varianza de la inflaci\'on; y (2) un componente impredecible $-\sigma\,dW(t)$, herencia de la estocasticidad del nivel de precios. En el caso determinista ($\sigma=0$), la ecuaci\'on de Fisher discreta se recupera: tasa real $\approx r - \mu$.
\end{intuicionmat}

\begin{intuicion}
Aun cuando el banco te paga una tasa cierta, tu poder adquisitivo real fluctua porque no sabes exactamente c\'omo va a cambiar el precio de los bienes. La parte aleatoria de la tasa real viene de los choques inesperados en la inflaci\'on, no del banco.
\end{intuicion}

%---------------------------------------------------------------
\section{Arbitraje y Completitud del Mercado\\
\textit{(Arbitrage and Market Completeness)}}
%---------------------------------------------------------------

La teor\'ia de precios de opciones descansa sobre dos nociones fundamentales: \textbf{mercados completos (complete markets)} y \textbf{ausencia de arbitraje (absence of arbitrage)}. Adem\'as, el modelo cl\'asico supone un \textbf{mercado perfecto (perfect market)}: sin costos de transacci\'on, sin restricciones al endeudamiento o a las ventas en corto, y sin limitaciones de ning\'un tipo sobre la selecci\'on de portafolios.

%---------------------------------------------------------------
\subsection{Noci\'on de Arbitraje \textit{(Notion of Arbitrage)}}
%---------------------------------------------------------------

\begin{intuicion}
Un arbitraje es, en esencia, un ``almuerzo gratis'': la posibilidad de ganar dinero sin arriesgar nada. En los mercados reales estas oportunidades existen brevemente y desaparecen r\'apido porque los inversores las aprovechan hasta que los precios se ajustan.
\end{intuicion}

Existe una \textbf{oportunidad de arbitraje (arbitrage opportunity)} cuando alguien puede obtener un retorno positivo con probabilidad positiva y sin ning\'un riesgo de p\'erdida. Matem\'aticamente, hay \textbf{arbitraje (arbitrage)} o \textbf{almuerzo gratis (free lunch)} en el mercado si existe una estrategia de negociaci\'on $\delta$ tal que:

\begin{equation}
  X^\delta(0) = 0, \quad X^\delta(T) \ge 0, \quad P\bigl[X^\delta(T) > 0\bigr] > 0 \quad \text{para alg\'un } T > 0
  \tag{3.42}
\end{equation}

\begin{intuicionmat}
La condici\'on (3.42) formaliza la idea de ``dinero gratis'': la estrategia no requiere inversi\'on inicial ($X^\delta(0) = 0$), nunca pierde ($X^\delta(T) \ge 0$), y existe al menos un escenario donde gana ($P[\cdots]>0$). Los modelos financieros est\'andar \textbf{asumen que no existen tales estrategias}, lo cual es el punto de partida para valorar activos y coberturas.
\end{intuicionmat}

donde $X^\delta(t)$ es la riqueza en el tiempo $t$ correspondiente a la estrategia $\delta$, y $P$ denota probabilidad.

\textbf{Nota importante:} la definici\'on de arbitraje es \textit{relativa al modelo}. Una estrategia compleja que requiere negociaci\'on continua constituye arbitraje solo si el modelo es correcto. En la pr\'actica, muchas ``oportunidades de arbitraje'' citadas por traders tienen baja probabilidad de p\'erdida, pero no cero, y por tanto no son arbitraje en sentido estricto.

%---------------------------------------------------------------
\subsection{Arbitraje en Modelos de Tiempo Discreto\\
\textit{(Arbitrage in Discrete-Time Models)}}
%---------------------------------------------------------------

\begin{intuicion}
En el modelo de un per\'iodo, el arbitraje se reduce a encontrar un portafolio de costo cero hoy que nunca pierda ma\~nana pero gane en al menos un estado. El ejemplo con t\'itulos Arrow-Debreu ilustra c\'omo la suma de sus precios debe igualarse al precio del activo libre de riesgo; de lo contrario, existe arbitraje.
\end{intuicion}

En el \textbf{modelo de un per\'iodo (single-period model)}, denotamos por $S_i^k(1)$ el pago del t\'itulo $i$ en el estado $k$ al tiempo $1$, $k=1,\ldots,K$. Existe una oportunidad de arbitraje cuando hay un portafolio $\delta$ tal que:

\begin{align}
  &\delta_0 B(0) + \delta_1 S_1(0) + \cdots + \delta_N S_N(0) = 0 \notag \\
  &\delta_0 B(1) + \delta_1 S_1^k(1) + \cdots + \delta_N S_N^k(1) \ge 0, \quad k = 1,\ldots,K
\end{align}

\begin{intuicionmat}
Las dos condiciones juntas dicen: (i) no se requiere capital inicial y (ii) el portafolio nunca pierde en ninguno de los $K$ estados posibles. Con al menos un estado con ganancia estricta, se confirma el arbitraje. Alternativamente, la condici\'on puede expresarse como tener costo inicial \textit{negativo} (se recibe dinero hoy) con pago nulo en todos los estados futuros.
\end{intuicionmat}

con al menos un estado $k$ tal que el pago sea estrictamente positivo. Una condici\'on equivalente es la existencia de un portafolio con costo negativo en $t=0$ y pago nulo en todos los estados en $t=1$:

\begin{align}
  &\delta_0 B(0) + \delta_1 S_1(0) + \cdots + \delta_N S_N(0) < 0 \notag \\
  &\delta_0 B(1) + \delta_1 S_1^k(1) + \cdots + \delta_N S_N^k(1) = 0, \quad k = 1,\ldots,K
\end{align}

\begin{intuicionmat}
Esta forma equivalente dice: recibes dinero hoy y no debes nada ma\~nana en ning\'un estado. Es la descripci\'on m\'as directa del ``almuerzo gratis''.
\end{intuicionmat}

\textbf{Ejemplo 3.5 (Arbitraje con T\'itulos Arrow-Debreu).} Considere un modelo de un per\'iodo con dos estados. Hay dos \textbf{t\'itulos Arrow-Debreu (Arrow-Debreu securities)} con precios $0.45$ cada uno, y un t\'itulo libre de riesgo que paga $1$ en cada estado y se vende a $0.95$. La estrategia de arbitraje consiste en comprar ambos t\'itulos Arrow-Debreu y vender en corto el t\'itulo libre de riesgo, generando un ingreso neto de $-0.45 - 0.45 + 0.95 = 0.05 > 0$ en $t=0$ y un pago nulo en $t=1$.

%---------------------------------------------------------------
\subsection{Arbitraje en Modelos de Tiempo Continuo\\
\textit{(Arbitrage in Continuous-Time Models)}}
%---------------------------------------------------------------

\begin{intuicion}
En tiempo continuo, la negociaci\'on continua puede generar estrategias ex\'oticas que producen arbitraje incluso en condiciones normales de mercado. Por eso se restringe el conjunto de portafolios a los ``admisibles''. Un caso concreto de arbitraje surge cuando dos activos con la misma volatilidad tienen distintos retornos esperados: se puede tomar una posici\'on larga/corta que genere ganancia sin inversi\'on inicial.
\end{intuicion}

Considera el modelo de \textbf{Merton-Black-Scholes (Merton-Black-Scholes model)} con dos acciones de igual volatilidad $\sigma$ conducidas por el mismo movimiento Browniano $W$. Existe arbitraje si $\mu_1 \ne \mu_2$. Suponiendo $\mu_1 > \mu_2$, la estrategia de vender en corto $S_1(0)/S_2(0)$ unidades de la acci\'on 2 y comprar una unidad de la acci\'on 1 tiene costo inicial cero y genera en $T$:

\begin{equation}
  S_1(T) - \frac{S_1(0)}{S_2(0)}S_2(T)
  = S_1(0)\,e^{\sigma W(T) - T\sigma^2/2}\!\left(e^{\mu_1 T} - e^{\mu_2 T}\right)
\end{equation}

\begin{intuicionmat}
La ganancia positiva se debe \'unicamente a la diferencia $e^{\mu_1 T} - e^{\mu_2 T} > 0$. Al tener la misma exposici\'on estoc\'astica $\sigma W(T)$, los riesgos se cancelan entre las dos posiciones y solo queda el diferencial de retornos esperados. Esto confirma que en un mercado libre de arbitraje, dos activos con \textit{id\'entica} estructura de riesgo deben tener el mismo retorno esperado.
\end{intuicionmat}

M\'as en general, para una econom\'ia con $d$ movimientos Brownianos independientes y $N$ t\'itulos riesgosos $S_i$, si existe una combinaci\'on lineal $C(t) = \sum_{i=1}^{N}\lambda_i S_i(t)$ con la misma estructura de volatilidad que otro t\'itulo $S$:

\begin{align}
  dC(t) &= \mu_1(t)C(t)\,dt + \sum_{j=1}^d \sigma_j(t)C(t)\,dW_j(t) \notag \\
  dS(t) &= \mu_2(t)S(t)\,dt + \sum_{j=1}^d \sigma_j(t)S(t)\,dW_j(t)
  \tag{3.43}
\end{align}

\begin{intuicionmat}
Las dos ecuaciones muestran que $C$ y $S$ comparten exactamente la misma exposici\'on a cada fuente de riesgo $dW_j(t)$, pero pueden diferir en su componente de deriva $\mu_1(t) \ne \mu_2(t)$. Esa diferencia es explotable: al mantener posiciones opuestas en $C$ y $S$, los t\'erminos estoc\'asticos se cancelan y s\'olo queda la diferencia de derivas, generando ganancia.
\end{intuicionmat}

Si existe un intervalo temporal donde $\mu_1(t) \ne \mu_2(t)$, se pueden construir estrategias de arbitraje tomando posiciones largas o cortas en $C$ y $S$.

%---------------------------------------------------------------
\subsection{Noci\'on de Mercados Completos \textit{(Notion of Complete Markets)}}
%---------------------------------------------------------------

\begin{intuicion}
Un mercado completo es aquel donde, con los activos disponibles, puedo construir cualquier flujo de pagos que imagine. En esos mercados, todo contrato tiene un precio \'unico y justo, lo que hace posible la valoraci\'on consistente de instrumentos complejos como las opciones.
\end{intuicion}

Un \textbf{contrato contingente (contingent claim)} es un contrato financiero con un pago aleatorio que puede ser positivo o negativo, dependiente de qu\'e estado del mundo se realice al vencimiento. Todos los t\'itulos descritos en el cap\'itulo son contratos contingentes.

Decimos que un portafolio \textbf{replica (replicates)} un contrato contingente cuando el pago del portafolio coincide con el pago del contrato en todos los estados posibles (con probabilidad uno). Se requiere que la estrategia replicante sea \textbf{autofinanciada (self-financing)}.

Un \textbf{mercado es completo (market is complete)} cuando es posible replicar cualquier contrato contingente con los t\'itulos existentes.

%---------------------------------------------------------------
\subsection{Mercados Completos en Tiempo Discreto\\
\textit{(Complete Markets in Discrete-Time Models)}}
%---------------------------------------------------------------

\begin{intuicion}
En el modelo binomial de un per\'iodo, con dos posibles estados futuros y dos instrumentos (un bono y una acci\'on), el sistema de ecuaciones para replicar tiene exactamente tantas ecuaciones como inc\'ognitas: siempre tiene soluci\'on \'unica. A\~nadir m\'as estados sin a\~nadir instrumentos deja el sistema sobredeterminado e irresoluble.
\end{intuicion}

Para el \textbf{modelo de un per\'iodo (single-period model)}, el valor terminal $X(1)$ de la estrategia replicante $\delta$ debe satisfacer:

\begin{equation}
  X(1) = \delta_0 B(1) + \delta_1 s_k = g(s_k), \quad k = 1,\ldots,K
  \tag{3.44}
\end{equation}

\begin{intuicionmat}
Este es un sistema de $K$ ecuaciones con 2 inc\'ognitas ($\delta_0$ y $\delta_1$). Solo tiene soluci\'on cuando $K \le 2$. El caso ``ideal'' es $K=2$ con sistema no degenerado: entonces se puede replicar cualquier contrato contingente y el mercado es \textbf{completo}. Si hay m\'as estados que instrumentos ($K > N$), el mercado es \textbf{incompleto}.
\end{intuicionmat}

\textbf{Ejemplo 3.7 (Replicaci\'on en modelo de un per\'iodo).} Con $r = 0.005$, $S(0) = 100$, $s_1 = 101$, $s_2 = 99$ y una opci\'on de compra europea $g[S(1)] = \max[S(1)-100,\,0]$, el sistema (3.44) es:
\begin{align*}
  \delta_0(1.005) + 101\delta_1 &= 1 \\
  \delta_0(1.005) + 99\delta_1  &= 0
\end{align*}
La soluci\'on es $\delta_0 = -49.2537$, $\delta_1 = 0.5$, con costo inicial $\delta_0 B(0) + \delta_1 S(0) = 0.746$, que es el precio justo de la opci\'on en este modelo.

\textbf{Condici\'on general (Teorema 3.3).} En el modelo Arrow-Debreu con $K$ estados y $N$ t\'itulos, sea $A$ la matriz $K \times N$ de pagos (columna $j$: pagos del t\'itulo $j$ en cada estado). Para replicar el vector de pagos $\vec{y} = (y^1,\ldots,y^K)$ se resuelve:

\begin{equation}
  \delta_1 x_1^i + \delta_2 x_2^i + \cdots + \delta_N x_N^i = y^i, \quad i = 1,\ldots,K
\end{equation}

\begin{intuicionmat}
El sistema tiene soluci\'on para \textit{cualquier} $\vec{y}$ si y s\'olo si $\text{rango}(A) = K$. Por lo tanto, la condici\'on necesaria y suficiente para completitud es que la matriz de pagos tenga rango m\'aximo (igual al n\'umero de estados). En particular, se necesitan al menos $K$ t\'itulos no redundantes.
\end{intuicionmat}

\textbf{Modelo binomial multiperi\'odo (CRR).} En el \textbf{modelo Cox-Ross-Rubinstein (Cox-Ross-Rubinstein model)}, el portafolio replicante se construye \textit{hacia atr\'as (backward induction)} desde el vencimiento. Para el nodo superior en el \'ultimo per\'iodo:

\begin{align}
  \delta_0^u(1)(1+r) + \delta_1^u(1)S^{uu} &= C^{uu} \notag \\
  \delta_0^u(1)(1+r) + \delta_1^u(1)S^{ud} &= C^{ud}
  \tag{3.45}
\end{align}

con soluci\'on:

\begin{align}
  \delta_1^u(1) &= \frac{C^{uu} - C^{ud}}{S^{uu} - S^{ud}} \notag \\
  \delta_0^u(1) &= \frac{1}{1+r}\!\left[C^{uu} - \delta_1^u(1)S^{uu}\right]
  \tag{3.46}
\end{align}

\begin{intuicionmat}
La ecuaci\'on para $\delta_1^u$ es el cociente entre el cambio en el valor del contrato y el cambio en el precio del subyacente. Este ratio se conoce como el \textbf{delta ($\Delta$)} del contrato: mide cu\'antas unidades del subyacente hay que mantener para replicar el contrato. Cambia en cada nodo del \'arbol.
\end{intuicionmat}

Repitiendo el proceso hacia atr\'as para $t=0$:

\begin{align}
  \delta_0(0)(1+r) + \delta_1(0)S^u &= C^u(1) \notag \\
  \delta_0(0)(1+r) + \delta_1(0)S^d &= C^d(1)
  \tag{3.47}
\end{align}

El n\'umero de acciones en el portafolio replicante es siempre de la forma:

\begin{equation}
  \delta_1 = \frac{\Delta C}{\Delta S}
  \tag{3.48}
\end{equation}

\begin{intuicionmat}
El \textbf{delta (delta)} $\delta_1 = \Delta C / \Delta S$ es el cociente de sensibilidades: cu\'anto cambia el precio del derivado cuando el subyacente cambia en una unidad. Es la pieza central de la cobertura din\'amica (\textit{dynamic hedging}) porque dice exactamente cu\'antas acciones mantener en cada instante.
\end{intuicionmat}

%---------------------------------------------------------------
\subsection{Mercados Completos en Tiempo Continuo\\
\textit{(Complete Markets in Continuous-Time Models)}}
%---------------------------------------------------------------

\begin{intuicion}
En tiempo continuo, el espacio de estados es infinito, lo que hace el an\'alisis m\'as sutil. Pero la idea es la misma: si hay tantos activos riesgosos como fuentes de incertidumbre (movimientos Brownianos), el mercado puede ser completo. La clave est\'a en la matriz de volatilidades: si tiene rango m\'aximo, todo riesgo puede ser gestionado.
\end{intuicion}

En los modelos de movimiento Browniano en tiempo continuo, un contrato contingente es cualquier contrato con pago terminal $C(T)$ que depende de la trayectoria de los movimientos Brownianos hasta $T$. El caso t\'ipico es la \textbf{opci\'on europea independiente de trayectoria (path-independent European option)} con pago $g(S(T))$.

Denotando $C(t)$ el precio en $t$ del contrato contingente, y asumiendo que todos los activos financieros satisfacen din\'amicas de la forma:

\begin{equation}
  dC(t) = a(t)\,dt + \sum_{j=1}^d b_j(t)\,dW_j(t)
  \tag{3.50}
\end{equation}

Para $N$ t\'itulos riesgosos con din\'amica:

\begin{equation}
  dS_i(t) = \mu_i(t)S(t)\,dt + \sum_{j=1}^d \sigma_{ij}(t)S_i(t)\,dW_j(t), \quad i=1,\ldots,N
  \tag{3.51}
\end{equation}

\begin{intuicionmat}
Cada t\'itulo $S_i$ tiene $d$ coeficientes de difusi\'on $\sigma_{i1},\ldots,\sigma_{id}$, uno por cada fuente de riesgo $W_j$. La matriz de volatilidades $\sigma(t)$ de tama\~no $N \times d$ resume toda la estructura de riesgo del mercado.
\end{intuicionmat}

La condici\'on de replicaci\'on $dX(t) = dC(t)$ se reduce a igualar los t\'erminos $dW_j$:

\begin{equation}
  \sum_{i=1}^N \pi_i(t)\,\sigma_{ij}(t) = b_j(t), \quad j = 1,\ldots,d
\end{equation}

\begin{intuicionmat}
Este es un sistema de $d$ ecuaciones con $N$ inc\'ognitas $\pi_i(t)$. Para replicar la exposici\'on estoc\'astica $b_j(t)$ del contrato contingente a la fuente de riesgo $j$, hay que combinar las exposiciones $\sigma_{ij}$ de los $N$ t\'itulos disponibles. Si el sistema tiene soluci\'on \'unica para cualquier vector $(b_1,\ldots,b_d)$, el mercado es completo.
\end{intuicionmat}

\textbf{Teorema 3.4 (Completitud en tiempo continuo).} Considere el modelo (3.51) con $d$ movimientos Brownianos independientes y $N$ t\'itulos riesgosos. Sea $\sigma(t)$ la matriz $N \times d$ de coeficientes de volatilidad. Una condici\'on necesaria y suficiente para que el mercado sea completo es que:

\begin{equation}
  \text{rango}[\sigma(t)] = d, \quad \forall\, t \in [0,T] \quad \text{(con probabilidad uno)}
\end{equation}

\begin{intuicionmat}
El rango de $\sigma(t)$ mide cu\'antas fuentes de riesgo pueden ser gestionadas. Si $\text{rango}(\sigma) = d$, cada una de las $d$ fuentes de incertidumbre tiene una contraparte en el mercado que permite replicarla. En particular, se necesita $N \ge d$: al menos tantos t\'itulos como movimientos Brownianos. El caso \'optimo es $N = d$ con $\sigma(t)$ no singular (invertible).
\end{intuicionmat}

Cuando $N = d$ y $\sigma(t)$ es no singular, el sistema de replicaci\'on tiene soluci\'on \'unica. Un t\'itulo es \textbf{redundante (redundant)} cuando su vector de coeficientes de difusi\'on es combinaci\'on lineal de los de otros t\'itulos.

\textbf{Ejemplo 3.9 (Replicaci\'on en el modelo Merton-Black-Scholes).} Con una acci\'on de coeficientes constantes $\mu$ y $\sigma$, y un bono con tasa $r$, si el contrato sigue:
\[
  dC(t) = a(t)\,dt + b(t)\,dW(t)
\]
la condici\'on de replicaci\'on exige:
\[
  \pi_0(t)\,r\,dt + \pi_1(t)[\mu\,dt + \sigma\,dW(t)] = a(t)\,dt + b(t)\,dW(t)
\]
lo que implica $\pi_1(t) = b(t)/\sigma$ y $\pi_0(t) = C(t) - \pi_1(t)$.

\textbf{Modelos incompletos.} Los modelos con \textbf{saltos (jumps)} en los precios o con \textbf{volatilidad estoc\'astica (stochastic volatility)} introducen fuentes de riesgo adicionales que generalmente no pueden ser replicadas con los activos disponibles, por lo que son modelos \textbf{incompletos (incomplete)}. Este libro se concentra en modelos completos por su mayor tractabilidad anal\'itica.

%---------------------------------------------------------------
\section{Ap\'endice: Demostraci\'on de la Regla de It\^o\\
\textit{(Appendix: Proof of It\^o's Rule)}}
%---------------------------------------------------------------

\subsection{Caso Simplificado de la Regla de It\^o \textit{(Simplified Case of It\^o's Rule)}}

\begin{intuicion}
La regla de It\^o es el an\'alogo estoc\'astico de la regla de la cadena del c\'alculo cl\'asico. La diferencia fundamental es que los procesos de Wiener tienen variaci\'on cuadr\'atica no nula, lo que introduce un t\'ermino extra de correcci\'on de segundo orden. Este t\'ermino ($\frac{1}{2}f''$) no existe en el c\'alculo ordinario y es la esencia del c\'alculo estoc\'astico.
\end{intuicion}

Se considera el caso m\'as simple de la \textbf{regla de It\^o (It\^o's rule)} (siguiendo la demostraci\'on de Steele, 2000):

\begin{equation}
  f(W(t)) - f(0) = \int_0^t f'(W(s))\,dW(s) + \frac{1}{2}\int_0^t f''(W(s))\,ds
  \tag{3.52}
\end{equation}

\begin{intuicionmat}
En el c\'alculo ordinario, $f(x(t)) - f(0) = \int_0^t f'(x(s))\,dx(s)$. La regla de It\^o agrega el t\'ermino $\frac{1}{2}\int f'' ds$, que surge porque el cuadrado del incremento Browniano $[dW(t)]^2 = dt$ (variaci\'on cuadr\'atica) no es despreciable. Esta correcci\'on es crucial en finanzas: por ejemplo, en el modelo Black-Scholes explica la aparici\'on del t\'ermino $-\frac{1}{2}\sigma^2$ en la soluci\'on expl\'icita del precio de la acci\'on.
\end{intuicionmat}

La demostraci\'on se basa en la expansi\'on de Taylor de $f$ y en el hecho de que, para el movimiento Browniano est\'andar $W$, la \textbf{variaci\'on cuadr\'atica (quadratic variation)} satisface $[W,W]_t = t$ (en el sentido $L^2$). Los t\'erminos de orden superior en la expansi\'on convergen a cero, mientras que el t\'ermino cuadr\'atico converge al integral $\frac{1}{2}\int_0^t f''(W(s))\,ds$.

%---------------------------------------------------------------
\section{Resumen y Conclusiones}
%---------------------------------------------------------------

\begin{enumerate}
  \item \textbf{Activos en tiempo continuo.} El activo libre de riesgo $B$ sigue $dB = Br\,dt$ y es localmente predecible. El nivel de precios $p$ sigue una din\'amica con componente Browniana y es impredecible. La tasa real de inter\'es resulta de deflactar $B$ por $p$, y hereda la impredecibilidad de la inflaci\'on.

  \item \textbf{Arbitraje.} Existe arbitraje cuando hay una estrategia con costo inicial cero, pago no negativo siempre y pago positivo con probabilidad estrictamente positiva. Los modelos est\'andar asumen ausencia de arbitraje como axioma fundamental.

  \item \textbf{Arbitraje en modelos discretos.} En el modelo de un per\'iodo, el arbitraje se detecta mediante un sistema lineal sobre los pagos del portafolio. En el modelo binomial CRR, la ausencia de arbitraje es verificable recursivamente.

  \item \textbf{Arbitraje en tiempo continuo.} Dos activos con la misma estructura de volatilidad pero distinta deriva generan arbitraje. La condici\'on de no-arbitraje exige igualdad de primas de riesgo entre activos con exposici\'on id\'entica.

  \item \textbf{Completitud.} Un mercado es completo si cualquier contrato contingente puede ser replicado. En tiempo discreto, la condici\'on es que la matriz de pagos $A$ tenga rango $K$ (tantas filas como estados). En tiempo continuo (Teorema 3.4), la condici\'on es que la matriz de volatilidades $\sigma(t)$ tenga rango $d$ (tantas columnas como fuentes de riesgo).

  \item \textbf{Delta.} El coeficiente de replicaci\'on del subyacente, $\delta_1 = \Delta C / \Delta S$, se denomina \textit{delta} del contrato. Es la pieza central de la cobertura din\'amica y var\'ia en cada nodo del \'arbol (discreto) o en cada instante (continuo).

  \item \textbf{Modelos incompletos.} Los modelos con volatilidad estoc\'astica o saltos suelen ser incompletos: existen riesgos que no pueden ser cubiertos con los activos disponibles. El libro se concentra en modelos completos por su mayor tractabilidad.
\end{enumerate}

\end{document}
